% !TeX root = ../main.tex
%------------------------------------------------------------------------------
%  Shafts, Caverns and Multiple Openings
%  Sources: Hoek & Brown (1980) Ch. 5; Brady & Brown (2006) Ch. 13;
%  Hoek & Moy on large powerhouse caverns; Barton et al. on large-span caverns.
%------------------------------------------------------------------------------
\chapter{Shafts, Caverns and Multiple Openings}
\label{ch:openings}

Running tunnels are the simplest underground structures and the least
representative of the works constructed. Shafts, portals, caverns,
junctions, cross-passages and station boxes carry a disproportionate share of
project risk: they are three-dimensional, they present several free surfaces
simultaneously, they are excavated in stages rather than continuously, and
they are the locations at which drives begin and end. Three themes recur
through the chapter. Geometry exerts a stronger influence than in a tunnel,
since shafts, caverns and junctions redistribute stress differently.
Excavation sequence affects the outcome, because the stress path to the final
geometry is not unique. Openings within one another's zone of influence
interact, so that none can be designed in isolation.

\section{Shafts and vertical openings}
\label{sec:shafts}

\subsection{Role and construction methods}
\label{sec:shaft-role-methods}

Shafts provide construction access, ventilation, drainage and, frequently, a
permanent operational function. Their cost per metre far exceeds that of
tunnelling, and their programme sits on the critical path because nothing
underground can start until the first shaft is complete.

Method follows ground and depth. In competent rock, conventional sinking by
drill and blast in short rounds with a sunk lining remains standard. In soft
ground the shaft must be supported before it is excavated, which leads to
diaphragm walls, secant or contiguous pile walls, or caisson sinking in which
the lining is built at the surface and sunk under its own weight as material is
excavated from within. Raise boring drills a pilot hole downward and reams
upward, which is efficient for ventilation and penstock shafts in rock but
requires the bore to stand unsupported between pilot completion and reaming,
and the pilot hole to be accurate over the full depth. Deviation that would
be trivial in a borehole becomes a serious problem when the hole is reamed to
a diameter of several metres. Shaft boring machines mechanise the full-face
operation.

Sinking is hazardous for geometric reasons. Everything
removed must be lifted through the single opening that people also use;
material dropped from the collar accelerates over the full depth; and water
entering the shaft collects at the working face rather than draining away.
For these reasons shaft sinking retains specialist crews and procedures
out of proportion to the volume excavated.

\subsection{Shaft stress and lining design}
\label{sec:shaft-stress-lining}

A vertical opening cuts the horizontal stress field, so the tangential stress
around a shaft is governed by the two horizontal principal stresses rather
than by the vertical stress. Where those differ, the stress concentrates on
the shaft wall in the direction perpendicular to the major horizontal stress.
In the high near-surface horizontal stresses of
Chapter~\ref{ch:in-situ-stress}, a shallow shaft can be more highly stressed
relative to strength than the deep tunnel it serves.

Spalling is therefore not confined to deep shafts, and its position around
the circumference
identifies the stress orientation (Figure~\ref{fig:shaft-stress}), so that observations during sinking
serve as an indirect stress measurement.

\begin{figure}[htbp]
\centering
\datafigure[0.62\textwidth]{ch24/shaft-stress}
\caption{Plan of a shaft with $\sigma_h = 0.5\,\sigma_H$. The tangential stress on the wall, plotted outward from the boundary, peaks at $3\sigma_H - \sigma_h$ perpendicular to $\sigma_H$, where spalling (carmine) forms, and falls to $3\sigma_h - \sigma_H$ in line with it.}
\label{fig:shaft-stress}
\end{figure}

A circular shaft lining works in hoop compression, which is an efficient load
path and permits a thin lining even at substantial water pressure. This efficiency
depends on circularity. An out-of-round lining develops bending, and
because the section was sized for compression it has little bending capacity;
construction tolerance is therefore a structural requirement rather than a
matter of appearance. Non-circular shafts lose the compression ring
and must be designed as framed or braced structures.

\subsection{Sinking in soft ground}
\label{sec:soft-ground-shafts}

In soft ground the governing checks are base stability and water. Base heave
occurs where the overburden pressure outside the shaft exceeds the bearing
capacity of the soil at the base; hydraulic uplift occurs where the water
pressure beneath a low-permeability layer exceeds the weight of the remaining
plug. Both are checked before excavation reaches the depth at which they
apply (Figure~\ref{fig:shaft-base-stability}), but they are not remedied by the same measures. Base heave is addressed
by deeper wall embedment, by a structural base slab, or by improving the soil
at founding level \citep{broms1967stability,davis1980stability}. Hydraulic
uplift is addressed by reducing the head beneath the plug, by taking the
cut-off into the aquitard, or by thickening the plug with jet grouting or
freezing (Chapter~\ref{ch:ground-improvement}). Dewatering inside the shaft
controls boiling at a granular base but does not reduce the head beneath a
low-permeability plug, and water removed from the shaft is weight no longer
resisting uplift.

\begin{figure}[htbp]
\centering
\datafigure[\textwidth]{ch24/shaft-base-stability}
\caption{The two base checks of a soft-ground shaft (schematic). Base heave: the soil outside, loaded by its overburden $\gamma z$, fails beneath the wall toe into the base at excavated depth $H$. Hydraulic uplift: the head $h_w$ in the aquifer acts beneath a plug of thickness $t$ whose weight $W$ resists it.}
\label{fig:shaft-base-stability}
\end{figure}

Where dewatering is unacceptable because of drawdown
(Chapter~\ref{ch:hydrogeology}), a cut-off wall taken into an aquitard, or a
grouted or frozen plug, becomes the only option, and the choice of shaft
construction method follows from it.

\section{Portals and access openings}
\label{sec:portals-openings}

\subsection{Portals}
\label{sec:portals}

A portal combines a slope problem with a tunnel problem, and failures
are disproportionately frequent there. The
cover is minimal, the ground is weathered, the stress is low so that
discontinuities are unclamped, and the first rounds are driven before any
support system is fully established.

The slope above and beside the portal must be stable before the tunnel is
started, and it must remain stable while the tunnel is driven beneath it. That
requires a kinematic assessment of the face and side slopes by the sliding
tests of Section~\ref{sec:mode-and-orientation}, and it usually requires support: soil
nailing, rock bolting, sprayed concrete facing, or a retaining structure.

The first rounds are then driven under the lowest cover on the project.
Pre-support is normal: an umbrella arch of grouted steel pipes installed from
the portal face carries the crown ahead of excavation
(Figure~\ref{fig:portal-section}), or forepoling spiles where the ground
only ravels, and the tunnel advances beneath it in short rounds with
immediate support. Both measures are described in
Section~\ref{sec:presupport-face}.
Where the ground is very poor the alternative is to build the first length as
cut-and-cover and backfill over it, which replaces a difficult tunnelling
problem with a conventional structural one.

\begin{figure}[htbp]
\centering
\datafigure[\textwidth]{ch24/portal-section}
\caption{Longitudinal section through a portal (schematic). The face slope is nailed before the first round, and the first rounds are driven in short lengths beneath overlapping stages of umbrella arch, under the lowest cover on the project and within the weathered zone.}
\label{fig:portal-section}
\end{figure}

\subsection{Openings, junctions and cross-passages}
\label{sec:junctions}

Forming an opening in a completed lining is a three-dimensional problem with
no closed-form solution. Removing part of a ring interrupts the hoop
compression that made it efficient, and load must redistribute around the
opening into the adjacent rings. The opening therefore requires a frame
(additional reinforcement, a steel or concrete surround, or specially designed
segments) installed before the breakout is made.

The ground at a junction is also more exposed. Two intersecting excavations
create additional free surfaces, and blocks that neither drive would release
on its own become removable at the intersection (Section~\ref{sec:block-and-wedge-construction}). Ground treatment before
breakout is common for this reason (Chapter~\ref{ch:ground-improvement}),
and cross-passages in soft ground
are frequently constructed under ground freezing because no other method
provides reliable support for a short opening between two completed tunnels.

The frame must be complete and the ground treated before the
lining is cut, and the design must state that sequence as a construction
requirement.

\subsection{Deep excavations and station boxes}
\label{sec:station-boxes}

A station box is a deep braced excavation, and its design is governed by wall
deflection rather than by collapse. Retaining wall selection (diaphragm wall,
secant pile, sheet pile) follows from depth, water and the movement that
adjacent assets will tolerate, and stiffer systems are chosen for movement
control rather than for strength.

Support may be by props, by ground anchors, or by top-down construction in
which the permanent floor slabs act as props as excavation proceeds beneath
them. Top-down is slower and costlier and gives the smallest movements, so it
is chosen where the surroundings are sensitive.

Base stability is checked as for a shaft, and the movement predictions feed
the damage assessment of Chapter~\ref{ch:settlement-damage}. Wall deflection
and ground settlement are related but not identical: settlement behind the
wall depends on the deflected shape as well as its magnitude, and a wall that
deflects mainly below excavation level produces a different surface trough
from one that cantilevers.

\section{Cavern layout and sequencing}
\label{sec:cavern-layout}

\subsection{Applications, orientation and span}
\label{sec:cavern-applications-orientation}

Large caverns house hydropower machine halls, metro station platforms, storage
for hydrocarbons and compressed air, and waste repositories. They
differ from tunnels in span, and span changes the nature of the problem
as well as its scale: the weight of a geometrically
similar wedge grows with the cube of the dimension
(Chapter~\ref{ch:structural-analysis}), the volume of overstressed rock grows,
and the probability of intersecting an adverse structure approaches certainty.

Shape follows the stress field. Mushroom, horseshoe and elliptical profiles
distribute boundary stress differently, and matching the profile to the stress
ratio reduces the peak concentration and the tensile zones
(Section~\ref{sec:elliptical}; Figure~\ref{fig:elliptical-stress}). The optimum is
elongation along the major principal stress in the ratio of the stresses, so
where the horizontal stress is high the profile should be wider than it is
tall, and where it is low a taller profile both lowers the peak concentration
and keeps the crown in compression.

\begin{figure}[htbp]
\centering
\datafigure[\textwidth]{ch24/elliptical-stress}
\caption{Boundary stress at the wall and crown of an elliptical opening of width $W$ and height $H$. The two are equal where $W/H = k$, the harmonic shape; a tall opening in high horizontal stress concentrates the stress in its crown.}
\label{fig:elliptical-stress}
\end{figure}

Orientation is the most valuable decision available and the most constrained.
The long axis should lie parallel to the major horizontal stress, so that the
cross-section is loaded by the smaller of the two horizontal stresses and the
peak boundary concentration, normally in the arch, is minimised
(Section~\ref{sec:elliptical}); the penalty is that the walls are then less
destressed than they would be with the major stress acting across the
section. The axis should also avoid running parallel to a
persistent discontinuity set to minimise wedge formation (Section~\ref{sec:construct-rotation}). These rarely
coincide, and the resolution requires identifying which mechanism governs:
stress in strong sparsely jointed rock at depth, structure in jointed rock
at moderate stress.

Empirical span guidance through the $Q$-system and equivalent dimension
extends to large caverns (Chapter~\ref{ch:empirical-design}), with an
excavation support ratio appropriate to a permanent civil structure, between
0.8 and 1.0 in Table~\ref{tab:esr}. It should be used with the reservation that the
case base thins rapidly with increasing span, and a chart extrapolated beyond
its data has no empirical support.

\subsection{Excavation sequencing}
\label{sec:cavern-sequencing}

A large cavern is excavated in stages. The standard sequence takes
a top heading first, installs the crown support from it, and then benches
downward in lifts, extending the sidewall support as each lift is exposed
(Figure~\ref{fig:cavern-sequence}).
Very large caverns may require multiple drifts, sidewall galleries or a pilot
tunnel through the crown.

\begin{figure}[htbp]
\centering
\datafigure[0.62\textwidth]{ch24/cavern-sequence}
\caption{A cavern of span $B = 24$~m and wall height 36~m excavated as a top heading (1) and four benches (2 to 5) (schematic). Crown and wall support alternate grouted bolts (malachite) with long cables (carmine), each row installed from the lift that exposes it.}
\label{fig:cavern-sequence}
\end{figure}

Sequencing serves two purposes. It limits the unsupported span at every stage,
so that no stage exceeds the ground's capacity even though the final span
does (Figure~\ref{fig:r3d-cavern-sequence}). It also provides progressive confirmation: the top heading exposes the
crown rock and its behaviour before the full geometry is committed, so the
design can be revised while revision is still cheap.

\begin{figure}[htbp]
\centering
\renderfigure{cavern_sequence}{%
  \rlab[west]{0.62,0.60}{0.83,0.70}{Cables}
  \rlab[west]{0.54,0.54}{0.83,0.52}{Bolts}
  \rlab[east]{0.40,0.60}{0.22,0.72}{Top heading}
  \rlab[east]{0.40,0.33}{0.22,0.30}{Next bench}
}
\caption{The cavern of Figure~\ref{fig:cavern-sequence} part-way through benching, with the rock removed from one side of the near end (schematic). The top heading and benches 2 and 3 are out and bench 4 (dark grey) is next; only the exposed lifts carry support, 14~m cables (carmine) alternating with 5~m bolts (malachite) in rings at an illustrative 3~m.}
\label{fig:r3d-cavern-sequence}
\end{figure}

The sequence is part of the design and belongs in the specification. Stress
redistributes at each stage, and the path to the final geometry affects both
the final stress state and the loads carried by support installed earlier.
A contractor who reorders the sequence for production reasons may
inadvertently change the design.

\subsection{Roof behaviour and reinforcement}
\label{sec:cavern-roof-reinforcement}

A cavern crown in bedded rock behaves as the voussoir arch of
Chapter~\ref{ch:ground-behaviour}, and at cavern span the maximum stable span
of an unreinforced bed is usually exceeded, so the roof must be reinforced
into a thicker composite beam or suspended from competent rock above.

Reinforcement at this scale relies on cables. Systematic bolting handles the
immediate skin and the smaller blocks; long pre-tensioned cables carry the
large wedges that extend beyond the reach of bolts. The two act together, and because
grouted bolts are stiffer than discretely anchored cables, the bolts load
first and the cables provide reserve at larger displacement.

Sidewalls are frequently more demanding than the crown. Tall walls in high
horizontal stress are destressed by excavation, which removes the clamping
that stabilises blocks, and steeply dipping foliation can buckle into the
opening. Wall support in large caverns is therefore designed separately
rather than as a continuation of crown support.

Observational control is the norm for cavern excavation. Instrumentation
installed from the top heading records crown displacement as benching
proceeds, trigger levels are tied to excavation stages, and support is
adjusted between lifts. This is the clearest routine application of the
observational method (Chapter~\ref{ch:movement-monitoring}), and it works because the sequence provides natural
decision points.

\section{Interaction between openings}
\label{sec:interaction}

\subsection{Zone of influence and stress shadows}
\label{sec:zone-influence-shadows}

An excavation disturbs the stress field over a region several diameters
across, conventionally bounded where the disturbance falls below about five
per cent of the field stress: for a circular opening in a hydrostatic
field, about one and three-quarter diameters clear of the boundary (Figure~\ref{fig:zone-of-influence}). Two excavations separated by more than that
do not interact; closer, they do, and each must be analysed with the other
present.

\begin{figure}[htbp]
\centering
\datafigure[0.62\textwidth]{ch24/zone-of-influence}
\caption{Stress about a circular opening of radius $a$ and diameter $D$ in a hydrostatic field $p_0$, against clear distance $d$ from the boundary. Both components depart from $p_0$ by $p_0(a/r)^2$, which falls to 5~per~cent (grey band) at $d = 1.74\,D$.}
\label{fig:zone-of-influence}
\end{figure}

Superposing the single-opening solutions is an approximation even in a wholly
elastic medium, because each solution leaves tractions on the boundary of the
other opening; it is acceptable while the separation keeps the zones of
influence clear of one another. Once yield zones from adjacent openings
overlap, superposition no longer applies and the problem must be solved as a
whole, numerically.

The disturbance is not uniformly adverse. An opening creates a destressed
region, a \emph{stress shadow}, in which a subsequent excavation experiences
lower stress than the far field. This can be exploited in sequencing, and it
carries the corresponding penalty that reduced stress means reduced
confinement, so a structurally controlled problem may be made worse by the
same shadow that relieves a stress-driven one.

\subsection{Pillar stress and strength}
\label{sec:pillar-stress-strength}

Where openings are separated by a pillar, the pillar carries the load that the
openings no longer support. The tributary area method quantifies this: the
pillar carries the full overburden over the plan area attributed to it, so
pillar stress rises as $1/(1-e)$ in the extraction ratio $e$, the proportion
of the plan area excavated (Figure~\ref{fig:pillar-stress}). It therefore rises very steeply once extraction is
high: pillar stress is twice the overburden stress at half extraction and
five times it at 80 per cent. The calculation sequence is set out in
Section~\ref{sec:seq-pillar} of Appendix~\ref{app:calculations}. The method assumes uniform extraction over a wide area and rigid
abutments, and it overestimates stress in a small group of openings where the
abutments carry a share.

\begin{figure}[htbp]
\centering
\datafigure[\textwidth]{ch24/pillar-stress}
\caption{Tributary area method. Left: square pillars of width $w$ separated by rooms of width $b$, each carrying the overburden on its dashed tributary area. Right: pillar stress against extraction ratio, twice the overburden stress at $e = 0.5$ and five times at 0.8; the square marks the layout drawn, $e = 0.61$.}
\label{fig:pillar-stress}
\end{figure}

Pillar strength differs from the strength of the rock of which it is made. It rises with
the width-to-height ratio, because a squat pillar develops confinement in its
core that a slender one cannot, and it falls with absolute size, because a
larger volume is more likely to contain a through-going weakness. Empirical
strength formulae~\citep{salamon1967pillar} capture both effects and are specific to the rock type and
the case base from which they were derived.

Factor of safety follows, with target values depending on duty and design
life. A temporary pillar in a mine and a permanent pillar between two metro
tunnels have different targets.

\subsection{Pillar failure and civil applications}
\label{sec:pillar-failure-civil}

Pillars fail progressively, and the mode of failure matters more than the
onset of yield. Spalling begins at the surface and works inward, reducing the effective
width and increasing the stress on what remains. A pillar with a high
width-to-height ratio can carry substantial load after yielding, because its
confined core survives; a slender pillar cannot, and its failure is sudden.

Whether that failure is stable or violent depends on relative stiffness. If
the surrounding rock mass is stiffer than the post-peak stiffness of the
failing pillar, the system sheds load in a controlled way; if it is softer,
the stored energy in the surrounding rock drives the failure and the pillar
bursts (Chapter~\ref{ch:high-stress}). Load transfer to neighbours on failure can
then trigger a cascade, and progressive collapse of a pillar system is a
well-documented mining failure mode.

Foundation failure is the remaining mode. A strong pillar bearing on a weak roof
or floor can punch into it, so that the system fails at a load well below the
pillar's own strength.

In civil works these principles apply to the rock or soil pillar between twin
running tunnels, to the walls between adjacent caverns, and to station
clusters where several openings sit within one another's zones of influence.
Minimum pillar width between twin tunnels is set by interaction rather than by
strength: too narrow, and the second drive's yield zone overlaps the first's,
increasing settlement and loading the first tunnel's lining. Construction
sequence then matters, because the second drive imposes additional load and
movement on a lining already built, which is why twin-tunnel settlement is
not twice the single-tunnel value
(Chapter~\ref{ch:settlement-damage}).

\chapterreferences
